\medskip
\noindent
13.3. \underline{The analytic Raynaud extension}.

We are now in a position to paraphrase, in analytic geometry, the
construction of the Raynaud extension given in 7. Suppose that $A$ is
as in 13.1, with $A=A^o$,
\footnote{The printed source has $A=A_o$ here. The required hypothesis,
and the one used below, is $A=A^o$; compare 13.1 and Lemma 13.2.1.}
and consider $R_o=R_s$, which is a free $\mathbb Z$-module of finite
type. After restricting $S$ to a sufficiently small open neighbourhood
of $s$, and interpreting $R_o$ as the group of germs of sections of $R$
near $s$, one obtains a homomorphism

\begin{equation}\tag{13.3.1}
 R_{oS}\longrightarrow R\subset E,
\end{equation}

\noindent
where $R_{oS}$ denotes the constant analytic group over $S$ with value
$R_o$. On choosing the neighbourhood of $s$ small enough that every
section of $\underline{\mathcal O}_S$ whose germ at $s$ is zero is itself
zero, one sees that the preceding homomorphism is injective. Its image is
therefore a constant analytic subgroup of $R$; it is plainly the largest
constant subgroup of $R$, and is called the \underline{fixed part} of
$R$ and denoted by $R^f$:

\begin{equation}\tag{13.3.2}
 R_{oS}\simeq R^f\hookrightarrow R,
 \qquad\text{inducing an isomorphism }R_s^f\simeq R_s.
\end{equation}

\noindent
We may then consider the quotient group

\begin{equation}\tag{13.3.3}
 A^\natural=E/R^f,
\end{equation}

\noindent
and, in view of (13.1.3), obtain a canonical exact sequence

\begin{equation}\tag{13.3.4}
 0\longrightarrow R/R^f\stackrel{i}{\longrightarrow}A^\natural
 \longrightarrow A\longrightarrow0.
\end{equation}

\noindent
Notice that, after restricting $S$ to a neighbourhood of $s$, $R^f$ is
always closed in $E$, and hence $A^\natural$ is separated over $S$.
Notice also that the homomorphism $A^\natural\longrightarrow A$ is
\'etale, and therefore induces an isomorphism

\begin{equation}\tag{13.3.5}
 \underline{\operatorname{Lie}}(A^\natural)
 \stackrel{\sim}{\longrightarrow}
 \underline{\operatorname{Lie}}(A)=\underline E.
\end{equation}

\noindent
Moreover, for every $n\geqslant0$, the canonical homomorphism
$A^\natural\longrightarrow A$ induces an isomorphism
$A_n^\natural\simeq A_n$ (since $R_n^f\stackrel{\sim}{\longrightarrow}
R_n$), and consequently an isomorphism

\begin{equation}\tag{13.3.6}
 \widehat{A^\natural}\stackrel{\sim}{\longrightarrow}\widehat A.
\end{equation}

\noindent
This last fact shows that $A^\natural$ does indeed play the role of the
group, likewise denoted by $A^\natural$, that we considered in 7. Here
it has been constructed without any hypothesis on the fibre
$A_s=A_o$ of the form ``an extension of an abeloid group by a torus''.
Such a hypothesis will nevertheless be necessary if we wish to recover
on $A^\natural$ the important extension structure (7.2.3). By contrast,
the homomorphism $A^\natural\longrightarrow A$ in (13.3.4) is essentially
transcendental in nature and has no counterpart in the purely algebraic
theory of \S~7. In \S~14 we shall return to the question of a purely
algebraic interpretation of the homomorphism
$i:R/R^f\longrightarrow A^\natural$ in (13.3.4).

Suppose first that we are given a complex subtorus

\begin{equation}\tag{13.3.7}
 T_o\subset A_o,
\end{equation}

\noindent
where by a ``\underline{complex torus}'' we mean an analytic group
isomorphic to a Cartesian power of $\mathbb C^*$. Let

\begin{equation}\tag{13.3.8}
 T=T_o\times S
\end{equation}

\noindent
be the constant torus over $S$ with value $T_o$. Proceeding as in 5.1,
one obtains a unique homomorphism (indeed, an immersion) of inductive
systems of analytic groups

\begin{equation*}
 \widehat T\longrightarrow\widehat{A^\natural},
\end{equation*}

\noindent
which over $S_o$ reduces to (13.3.7). By 13.2.2, this homomorphism comes
from a unique germ of homomorphisms

\begin{equation}\tag{13.3.9}
 T\longrightarrow A^\natural,
\end{equation}

\noindent
and, after taking $S$ small enough around $s$, we may suppose that it is
defined over all of $S$. The corresponding homomorphism

\begin{equation*}
 \underline{\operatorname{Lie}}(T)\longrightarrow
 \underline{\operatorname{Lie}}(A^\natural)\simeq\underline E
\end{equation*}

\noindent
is injective over each $S_n$, hence over $\widehat S$, and therefore over
$\underline{\mathcal O}_s$ by flat descent. Thus there is a neighbourhood
of $s$ in $S$, which we may take to be all of $S$, on which
$\underline{\operatorname{Lie}}(T)$ identifies with a locally direct
summand $\underline E^t$ of $\underline E$; denote the corresponding
vector subbundle of $E$ by $E^t$:

\begin{equation}\tag{13.3.10}
 \underline{\operatorname{Lie}}(T)\xrightarrow{\ \sim\ }
 \underline E^t\subset\underline E\simeq
 \underline{\operatorname{Lie}}(A^\natural)\simeq
 \underline{\operatorname{Lie}}(A),
 \qquad E^t\subset E.
\end{equation}

\noindent
Likewise, let

\begin{equation}\tag{13.3.11}
 R_o^t\subset R_o,
 \qquad R^t=R_{oS}^t\subset R_{oS}\simeq R^f
\end{equation}

\noindent
denote respectively the sublattice of $R_o$ defined by the subtorus
$T_o$ of $A_o$, and the corresponding constant analytic group over $S$.
Then the homomorphism (13.3.9) corresponds simply to the homomorphism

\begin{equation*}
 E^t/R_{oS}^t=E^t/R^t\longrightarrow E/R_{oS}=E/R^f
\end{equation*}

\noindent
induced by the inclusions

\begin{equation*}
 E^t\hookrightarrow E,
 \qquad R^t\hookrightarrow R^f.
\end{equation*}

\noindent
Using the fact that (13.3.9) induces an immersion on the fibre over $s$,
we conclude that the homomorphism

\begin{equation*}
 R_o/R_o^t\longrightarrow E_o/E_o^t=(E/E^t)_o
\end{equation*}

\noindent
is injective. It follows at once that

\begin{equation*}
 R_o/R_o^t\longrightarrow(E/E^t)_{s'}
\end{equation*}

\noindent
is still injective for $s'$ near $s$. Thus, after shrinking $S$ around
$s$, we may suppose that (13.3.9) \underline{is a closed immersion}.

We are therefore led to introduce

\begin{equation}\tag{13.3.12}
 B=A^\natural/T\simeq(E/E^t)/(R^f/R^t),
\end{equation}

\noindent
so that $A^\natural$ acquires an extension structure

\begin{equation}\tag{13.3.13}
 0\longrightarrow T\longrightarrow A^\natural
 \longrightarrow B\longrightarrow0,
\end{equation}

\noindent
where $B$ is a relative analytic group over $S$, smooth and separated
over $S$, with connected fibres. By construction we therefore have

\begin{equation}\tag{13.3.14}
 \underline{\operatorname{Lie}}(B)\simeq
 \underline E/\underline E^t,
\end{equation}

\noindent
and the exponential exact sequence for $B$ is

\begin{equation}\tag{13.3.15}
 0\longrightarrow R^f/R^t\longrightarrow E/E^t
 \longrightarrow B\longrightarrow0.
\end{equation}

\noindent
Notice that

\begin{equation}\tag{13.3.16}
 R^f/R^t=(R_o/R_o^t)_S
\end{equation}

\noindent
is a constant analytic group over $S$. It follows from 13.1.5 that $B$
is an abeloid group over $S$ in a neighbourhood of $s$ if and only if
its fibre at $s$,

\begin{equation*}
 B_o=A_o/T_o,
\end{equation*}

\noindent
is an abeloid analytic group; in that case $A_o$ appears as an extension
of such a group $B_o$ by a torus $T_o$. \underline{Every way of writing
$A_o$ as an extension of an abeloid group by a complex torus $(*)$
therefore gives rise on $A$, in a neighbourhood of $s$, to a unique
extension structure} (13.3.13) \underline{of an abeloid group $B$ by a
torus $T$, inducing the given extension structure at $s$}.
